\documentclass[12pt, letterpaper]{article}
\usepackage[margin=1.2in]{geometry}
\usepackage{ebgaramond}
\usepackage[T1]{fontenc}
\usepackage{microtype}
\usepackage{setspace}
\usepackage{parskip}
\usepackage{titling}
\usepackage{fancyhdr}

\setlength{\parindent}{2em}
\onehalfspacing

\pagestyle{fancy}
\fancyhf{}
\fancyhead[L]{\small\textit{The One That Didn't}}
\fancyhead[R]{\small\textit{NASA Mission Series v1.11}}
\fancyfoot[C]{\small\thepage}
\renewcommand{\headrulewidth}{0.4pt}

\title{\Huge\textbf{The One That Didn't}\\[0.5em]\large\textit{A short story from Mission 1.11 --- Explorer 5}}
\author{NASA Mission Series\\[0.3em]\small Miles Tiberius Foxglove --- Mukilteo, WA}
\date{March 30, 2026}

\begin{document}
\maketitle
\thispagestyle{empty}

\vspace{1em}
\noindent\rule{\textwidth}{0.4pt}
\vspace{1em}

\noindent The payload knew nothing about itself. It was not aware, not sentient, not anything that thinks or hopes or understands the difference between up and down. It was forty pounds of instruments in a sealed cylinder, wired to a transmitter, bolted to a cluster of solid rocket motors, and it did not know it was going to space. It did not know it was going to fail.

But if it had known --- if the Geiger tubes and the magnetometer and the crystal detectors had been capable of something beyond counting and measuring and transmitting numbers to a ground station in Florida --- here is what it would have known:

It was the fifth. Explorer 1 had gone up in January and found the radiation belts that nobody expected. Explorer 2 had failed --- fourth stage didn't ignite. Explorer 3 had succeeded, confirming what Explorer 1 discovered. Explorer 4 had gone up in July with scintillation counters designed for the Argus nuclear tests. And now, August 24, 1958, Explorer 5 sat atop a Juno I rocket at Cape Canaveral, carrying improved instruments to map the Van Allen belts with better precision than any of its predecessors.

The instruments were ready. The rocket was fueled. The spinning tub --- the cylindrical assembly that held the upper stages --- was spun up to 450 revolutions per minute, a blur of angular momentum that would keep the trajectory stable after separation from the booster. Everything worked. Everything was nominal. The countdown reached zero and the Rocketdyne engine ignited and the rocket climbed into the Florida night sky on a column of fire that was visible from Cocoa Beach.

\begin{center}
$\bullet\quad\bullet\quad\bullet$
\end{center}

In Jorge Luis Borges's Library of Babel, there exists a book for every possible combination of characters. Most books are gibberish. Some contain a single coherent sentence surrounded by chaos. A very few contain truth. Somewhere in that library, on a shelf in a hexagonal room lit by a dim lamp that never goes out, there is a book that contains the complete dataset that Explorer 5 would have collected if it had reached orbit. Every magnetometer reading, every Geiger count, every temperature measurement, every telemetry frame --- all of it written in a volume that no one will ever read, because the experiment never ran.

The library contains the book. The universe does not contain the data.

\begin{center}
$\bullet\quad\bullet\quad\bullet$
\end{center}

One hundred and fifty seconds into the flight, the first stage burned out. The Rocketdyne engine fell silent. The Juno I was coasting now, rising on momentum alone, the spinning tub humming at 450 RPM above the spent booster. The separation mechanism armed. Pyrotechnic bolts fired. Springs pushed. The spinning cluster began to move away from the booster body.

The clearance between the edge of the spinning tub and the inner wall of the booster was measured in centimeters. Not many centimeters. The design called for five. Manufacturing tolerance consumed one. Thermal expansion from the engine burn consumed another. A slight mass eccentricity in the spinning cluster --- unavoidable in a system balanced to the limits of 1958 technology --- consumed a third. What remained was a margin of two centimeters between mission success and mission failure, and on August 24, 1958, two centimeters was not enough.

The edge of the spinning tub touched the booster.

Not a collision. Not a crash. A brush --- the lightest possible contact between two pieces of aluminum moving at slightly different velocities, one spinning, one not. The contact lasted perhaps ten milliseconds. The force was modest. The damage to the hardware was negligible. But the impulse was off-axis, and an off-axis impulse applied to a spinning body does something that no amount of engineering can undo: it induces nutation.

Nutation is the wobble of a spinning top that has been bumped. The spin axis, which should point straight ahead, begins to trace a cone. The angle of the cone is the nutation angle, and once it starts, it cannot be stopped --- not in a system with no active attitude control, which the Juno I upper stages did not have. The gyroscopic stability that was supposed to keep the trajectory true was now keeping the wobble true instead. The angular momentum was conserved. The direction was wrong.

The upper stages fired on schedule. The solid rocket motors ignited and burned with the same thrust and the same duration as they would have in a nominal flight. But the thrust vector was tilted by the nutation angle, which meant the velocity increment was pointed in the wrong direction. The payload gained speed. It gained altitude. For a few minutes, it was higher than any human artifact had been since Explorer 4 a month before. But it was not going fast enough in the right direction, and it was going too fast in the wrong direction, and the difference between those two things is the difference between orbit and not-orbit.

\begin{center}
$\bullet\quad\bullet\quad\bullet$
\end{center}

The payload reached apogee. For a moment --- a fraction of a second that the instruments could have measured if the mission planners had thought to include an accelerometer, which they hadn't, because you don't need an accelerometer when you're going into orbit --- the payload was weightless. It was in freefall. It was, briefly, a satellite of the Earth, tracing the beginning of an ellipse that, if completed, would have carried it around the planet in approximately ninety minutes.

The ellipse was not completed. The perigee of the orbit --- the lowest point --- was inside the atmosphere. Not by a lot. Not by a hundred kilometers. By enough. The payload began to descend. The transmitter was still transmitting. The instruments were still measuring. The Geiger tubes were counting cosmic rays and Van Allen belt particles, and the magnetometer was recording the magnetic field, and the data was being telemetered to the ground station, and for a few minutes the mission was doing exactly what it was designed to do, except that it was doing it on a trajectory that would end in the Atlantic Ocean instead of continuing around the Earth.

Borges wrote about the Aleph --- a point in space that contains all other points, where everything that happens is visible simultaneously. The payload, for a few minutes at apogee, was an Aleph of measurement: it could see the radiation belts from inside them, could count the particles, could feel the magnetic field. It was inside the thing it was built to study. And then it fell out.

\begin{center}
$\bullet\quad\bullet\quad\bullet$
\end{center}

In the forests of the Pacific Northwest, in the damp shade beneath Sitka spruce and western hemlock where the light filters through a hundred meters of canopy, \textit{Amanita muscaria} pushes through the forest floor. The mushroom begins as an egg --- a small white sphere enclosed in a universal veil, a membrane that wraps the entire developing fruiting body in a protective skin. As the mushroom grows, the veil stretches. As the mushroom pushes upward, the veil tears. The stipe elongates, the cap expands, and the membrane that protected it breaks apart into patches --- white fragments that remain on the red cap like snow on a mountain.

The veil-breaking is the mushroom's separation event. The protective enclosure opens, the organism emerges, and in that moment of transition --- between contained and free, between protected and exposed --- everything depends on clearance. If the veil tears cleanly, the cap emerges smooth and symmetrical, the white patches distributed evenly, the stipe straight and strong. If the veil catches, if something snags, if the clearance between the expanding cap and the tearing membrane is not quite sufficient, the cap deforms. The patches pile up on one side. The stipe bends. The mushroom grows, but it grows wrong, and no amount of subsequent growth can correct the error introduced at the moment of emergence.

Explorer 5's separation was its veil-breaking. The payload was enclosed in the launch vehicle, protected during ascent. At the moment of emergence --- when the springs fired and the tub began to separate --- the clearance was not sufficient. The edge caught. The trajectory deformed. The mushroom grew, but it grew wrong.

\begin{center}
$\bullet\quad\bullet\quad\bullet$
\end{center}

The ground station at Cape Canaveral tracked the signal for approximately four minutes after separation. The telemetry showed nominal instrument operation --- all channels reporting, all measurements within expected ranges. The data was good. The instruments worked. The Geiger tube counted particles. The magnetometer measured field strength. For four minutes, Explorer 5 was a functioning scientific satellite producing exactly the data it was designed to produce.

Then the signal began to fluctuate. The tumbling payload rotated its antenna away from the ground station, then back, then away again, each rotation reducing the received signal strength. The S-meter at the tracking station swung from S7 to S2 to S7 to S1, the carrier cutting in and out as the antenna pattern swept through the ground station's beam. The operator logged the anomaly. The flight director noted the unexpected rotation. The trajectory team plotted the orbit determination and found no orbit to determine.

The signal faded. The static rose. The carrier dropped below the noise floor. The operator turned the receiver gain to maximum and heard nothing but the hiss of cosmic background and the sixty-cycle hum of the ground equipment. He wrote in the log: ``Signal lost. No acquisition on subsequent pass.'' There would be no subsequent pass. Explorer 5 was in the Atlantic Ocean.

\begin{center}
$\bullet\quad\bullet\quad\bullet$
\end{center}

Borges imagined a garden of forking paths where every possible outcome of every possible event exists simultaneously, each fork creating a new universe in which that particular outcome is the one that happened. At the moment of separation, the garden forked. In one path, the tub cleared the booster by two centimeters, the nutation angle was zero, the upper stages fired true, and Explorer 5 entered a 300-by-2200-kilometer orbit with an inclination of 33 degrees and a period of 107 minutes. In that path, the instruments mapped the Van Allen belts with unprecedented precision, discovered the slot region between the inner and outer belts three years before it was actually discovered, and transmitted data for four months until the battery died.

In the other path --- our path --- the tub touched the booster.

Both paths exist in the library. Both books are on the shelves. The book of Explorer 5's actual data --- four minutes of telemetry from a suborbital trajectory --- is thin and incomplete and ends mid-sentence. The book of Explorer 5's intended data --- four months of orbital measurements --- is thick and complete and ends with the battery voltage dropping below the transmitter threshold and the signal fading gracefully into the noise, the way a satellite is supposed to die.

Nobody will ever read the second book. But it exists. The library is complete.

\begin{center}
$\bullet\quad\bullet\quad\bullet$
\end{center}

The \textit{Amanita} is still growing. Under the Sitka spruce, in the rain-soaked duff of the Pacific Northwest forest floor, a fly agaric pushes its red cap through the broken veil, white patches stark against the crimson surface. The mushroom does not know about Explorer 5 or Juno I or the Van Allen belts. It does not know that separation is the most dangerous moment. It knows only this: the veil must break for the spores to be released, and the breaking is never clean, and the patches that remain are the evidence of the breaking, and the evidence is beautiful.

Two centimeters. The distance between orbit and ocean, between data and possibility, between the book that was written and the book that was not. The distance between the edge of the spinning tub and the wall of the booster, measured in the currency of 1958 manufacturing tolerance, which was insufficient for the question being asked. The distance between the cap of the mushroom and the tearing veil, which is always insufficient, which is why the patches are always there, which is why the \textit{Amanita} is always imperfect, which is why it is always beautiful.

The one that didn't. Not the one that failed --- failure implies an ending, and Explorer 5 did not end, it simply took the other path. The instruments are in the Atlantic. The data is in the library. The mushroom is in the forest. The garden forks, and forks, and forks, and somewhere down every path there is a version of August 24, 1958, where two centimeters was enough.

\end{document}
